\documentclass[11pt]{article}
\usepackage[utf8]{inputenc}
\usepackage[T1]{fontenc}
\usepackage[spanish]{babel}
\usepackage{geometry}
\geometry{a4paper, margin=2.5cm}
\usepackage{hyperref}
\usepackage{listings}
\usepackage{courier}
\lstset{basicstyle=\ttfamily\small,breaklines=true}

\begin{document}
\title{Manual Técnico - Quibdó Seguro}
\author{Equipo de Desarrollo}
\date{2026}
\maketitle

\section{Arquitectura}
La aplicación está organizada como una UI monolítica (Laravel) y varios microservicios:
\begin{itemize}
  \item Auth Service (puerto 8001)
  \item Incident Service (8002)
  \item Rewards Service (8003)
  \item Notification Service (8004)
  \item Analytics Service (8005)
  \item API Gateway (8000)
\end{itemize}

\section{Variables de entorno críticas}
\begin{itemize}
  \item `GATEWAY_URL`
  \item `MONGODB_URI`, `DB_DATABASE`
  \item `SESSION_DRIVER`, `LOG_CHANNEL`
  \item `SANCTUM` y modelo `personal_access_token_model` en `auth-service`
\end{itemize}

\section{Endpoints importantes}
\begin{lstlisting}
POST /auth/register    (gateway -> auth-service)
POST /auth/login
GET  /auth/profile
POST /auth/logout
POST /incidentes
GET  /stats
\end{lstlisting}

\section{Docker Compose (ejemplo desarrollo)}
Archivo `docker-compose.dev.yml` (ejemplo mínimo):
\begin{lstlisting}
version: '3.8'
services:
  mongo:
    image: mongo:6
    ports: ["27017:27017"]
    volumes:
      - mongo-data:/data/db

  gateway:
    build: ./microservices/gateway
    ports: ["8000:8000"]
    environment:
      - MONGODB_URI=mongodb://mongo:27017/quibdo_seguro

  auth-service:
    build: ./microservices/auth-service
    ports: ["8001:8001"]
    environment:
      - MONGODB_URI=mongodb://mongo:27017/quibdo_seguro

volumes:
  mongo-data:

\end{lstlisting}

\section{Hosting gratuito sugerido}
\begin{itemize}
  \item Documentación estática (PDF / Markdown): GitHub Pages o Netlify (gratis)
  \item Aplicación completa: Render / Railway (planes gratuitos limitados) o Heroku (si disponible)
  \item Para pruebas, usar `docker-compose` + Cloudflare Tunnel para exponer localmente en un dominio temporal
\end{itemize}

\section{Accesibilidad y diseño}
Usar variables CSS con contraste adecuado. Tipografías y tamaños legibles para usuarios mayores.

\section{Pruebas recomendadas}
Unit tests (PHPUnit), pruebas E2E (cypress o Dusk), contract tests OpenAPI.

\end{document}
